\section{Overview of computing infrastructures}

A computing infrastructure is a technological foundation that provides the necessary hardware and software for computation to other systems and devices. Modern infrastructures are composed of varied, interconnected components.

\subsection{The computing continuum}

The computing continuum is a model describing a distributed environment that seamlessly incorporates processing elements across endpoints, the edge and the cloud.

The choice of using this paradigm is driven by two main factors. First, the \textbf{exponential growth in data generated} by connected devices requires innovative solutions for efficient processing and analysis. Second, traditional centralized data centers struggle with scalability and latency (real-time processing), hindering the performance in resource-intensive applications.

The continuum spans a massive spectrum of hardware, generally categorized into four operational tiers:
\begin{enumerate}
    \item \textbf{Internet of Things (IoT) endpoints} - They are crucial for real-time data collection in distributed environments. An IoT device is an everyday object embedded with sensors, software, and internet connectivity. They offer highly pervasive wireless connections and very low costs. However, they suffer from low computing ability, tight battery constraints, and minimal memory.
    \item \textbf{Embedded devices} - They bridge the gap between microcontrollers and traditional servers. An embedded PC can be programmed like standard PCs and provides a higher performance computing unit. However, they feature pretty high power consumption and often require custom hardware design efforts.
    \item \textbf{Edge and Fog computing systems} - They utilize local servers to process data near its source. The key difference between them is the location of the processing. Edge computing places the intelligence directly in the device, while fog computing moves the processing to intermediary device linked to a local area network. These computing systems provide high computational capacity, enable distributed computing, and drastically reduce latency in decision-making. However, they necessitate a constant power connection and rely on a persistent connection with the cloud.
    \item \textbf{Data centers} - They are at the top of the continuum, and provide the remote, large-scale processing capabilities essential for handling vast data volumes. Data centers lower the IT costs while delivering high-performance, universal data access, and seemingly unlimited storage capacity with increased data reliability. However, they require a constant high-speed internet connection, consume massive amounts of power, and introduce latency when taking decisions from distant endpoints.
\end{enumerate}
